\section{问题重述}

\subsection{研究背景}

微网通过协调光伏、储能与外网购电满足本地负载，可利用低价购电、富余光伏充电及高价时段放电降低运行费用\cite{lasseter2002,katiraei2008}。由于负载、光伏及电价随时间变化，购电计划还需应对预测偏差引起的紧急补购与调整费用。本文利用题给数据，在满足供需平衡和储能运行约束的前提下，研究不同信息条件下的购电与充放电策略。

\subsection{问题提出}

\textbf{问题一：}在给定典型日电价、负载与光伏预测的条件下，制定全天购电及储能充放电计划，使日初、日末储电量相同，并最小化购电费用。

\textbf{问题二：}在负载与光伏未知、每日分时电价相同的条件下，每天零点制定当天购电计划；结合实际运行中的紧急补购，给出指定日期及全年购电结果。

\textbf{问题三：}利用每天 $0/6/12/18$ 点发布的光伏预报制定并更新购电策略，综合计划、调整与紧急购电费用，判断引入日内调整是否有必要。

\textbf{问题四：}在波动电价条件下重新求解问题二、三，分别给出日前购电与日内调整策略的计算结果。各问的模型衔接见图~\ref{fig:overall_framework}。
